\begin{abstract}
大喜利投稿サイト「ボケて」のお題画像の約6割には読める文字があり，その文字が回答者への題であるか，画像の素材の一部であるかによって，ボケの生成に必要な処理が変わる．
本稿では，約682万件のお題画像へのラベルの付与において，画像中の文字をどのように定義し，その定義がラベルの再現性，教師となる視覚言語モデル (Vision-Language Model，VLM) と生徒となる分類器の精度，全データのラベルの分布にどのような影響を与えたかを報告する．
従来の定義では，お題文を「投稿者が回答者に向けて書き加えた文」とし，元の素材の字幕や看板を除いていた．
この定義の下では，誤りの多くがテロップや帯の文字の出所の判断に集中し，VLM は投稿者が加えた文と元の字幕を区別できなかった．
そこで，文字を写っている場所による4つの層と，文の形による5つの形に分け，お題文の有無を属性から規則で導く定義を提案した．
1,620枚の目視ラベルで評価した結果，提案の属性は付け直しで記号の94\%が完全に一致し，選んだ規則で導いたお題文の有無は従来の定義と全データの推定で98.5\%（Cohen の $ \kappa $ 0.90）一致した．
また，規則を変えるだけで，お題文と判定される画像の割合が7.4\%から16.7\%まで変わることを示した．
一方，Qwen3-VL-32B は層を区別できず，教師としての F1 は0.43にとどまり，1,620枚で学習した生徒も既存のパッケージの F1 0.83 に届かなかった．
\par
\noindent
{\bf キーワード} 大喜利，データアノテーション，ラベル設計，視覚言語モデル，画像中の文字
\end{abstract}

\section*{Abstract}
About 60\% of the prompt images on the Japanese Ogiri site Bokete contain readable text, and whether the text is a prompt addressed to the answerers or part of the source material changes the processing needed to generate a witty response.
This paper reports how we defined text in prompt images when annotating about 6.82 million images, and how the definition affected label reproducibility, the accuracy of a vision-language model (VLM) teacher and a student classifier, and the label distribution of the whole dataset.
The previous definition treated a prompt sentence as text added by the poster for the answerers and excluded captions and signs in the source material.
Under this definition, most errors concentrated on judging the origin of captions and bands, and the VLM could not distinguish text added by the poster from original subtitles.
We therefore propose to split text into four layers by where it is placed and five forms by its sentence form, and to derive the presence of a prompt sentence from these attributes by a rule.
On 1,620 manually labeled images, 94\% of the attribute codes were reproduced exactly on relabeling, and the presence of a prompt sentence derived by the selected rule agreed with the previous definition on 98.5\% of the population (Cohen's $ \kappa $ 0.90).
We also show that changing the rule alone moves the share of images with a prompt sentence from 7.4\% to 16.7\%.
However, Qwen3-VL-32B could not distinguish the layers and reached an F1 of only 0.43 as a teacher, and students trained on the 1,620 images did not reach the F1 of 0.83 of the existing package.
\par
\noindent
{\bf Key words} Ogiri, Data Annotation, Label Design, Vision-Language Model, Text in Images
